%%=============================================================================
%% Voorwoord
%%=============================================================================

\chapter*{\IfLanguageName{dutch}{Woord vooraf}{Preface}}%
\label{ch:voorwoord}

%% TODO:
%% Het voorwoord is het enige deel van de bachelorproef waar je vanuit je
%% eigen standpunt (``ik-vorm'') mag schrijven. Je kan hier bv. motiveren
%% waarom jij het onderwerp wil bespreken.
%% Vergeet ook niet te bedanken wie je geholpen/gesteund/... heeft

Ik heb dit onderwerp gekozen uit de lijst van bachelorproeven van meneer Clauwaert. Origineel sprak het onderwerp me aan omdat ik het antwoord ook wilde weten op de vraag. Toen was ik ook veel bezig buiten school met Linux en leek me de perfecte mogelijkheid om me te verdiepen in het onderwerp. Buiten deze bachelorproef zou ik me niet zo verdiept hebben in rootkits. Tijdens het uittesten van de rootkits realiseerde ik me dat ik niet eens de basis bezitten om de concepten te begrijpen. Daarom krijg ik ook motivatie voor een demonstratie en uit te leggen hoe rootkits werken en hoe te beschermen ertegen. Ik wil graag Dhr. T. Clauwaert en Dhr. B. Van Vreckem bedankt voor het bijsturen van mijn doel voor deze bachelorproef. 